\section{Troisième et Quatrième semaines : Animation de l'atelier}

  Durant les deux dernières semaines, nous avons eu l'occasion d'animer l'atelier auprès des élèves de seconde.
Nous avons animé deux ateliers par jour : j’étais chargé de raconter l’histoire, tandis que mon binôme guidait les élèves à travers les énigmes. Certaines énigmes ont parfois été omises pour certains groupes, au profit d'aborder plus profondément certaines notions, ou bien de privilégier une partie quand celle-ci suscitait plus d'entrain pour un groupe. \par
  L'animation de l'atelier a été une expérience très enrichissante, car elle m'a permis de mettre en pratique les connaissances acquises lors de la phase de conceptualisation et de concrétisation de l'atelier.
  \subsection{Réactions des élèves}
  Les élèves ont été très réceptifs à l'atelier et ont montré un grand intérêt pour les activités proposées. Ils ont posé de nombreuses questions et ont été très curieux de comprendre le fonctionnement du modèle de langage.  \par 
  Nous avons toutefois pu constater que les élèves avaient globalement un niveau de connaissances très limité, d'abord bien sûr d'un point de vue théorique, mais aussi surtout d'un point de vue de l'utilisation et des limites de tels outils. En effet, il arrive à certains élèves de croire aveuglément aux réponses des modèles qu'ils utilisent au quotidien, et de totalement l'humaniser. Nous avons alors martelé que ces modèles sont avant tout mécaniques, que ceux-là obéissent à des lois mathématiques, et que leurs réponses ne sont fiables qu'en apparence, qu'elles peuvent être biaisées par la formulation de la question, et qu'elles ne sont pas toujours objectives. \par
 Je suis déçu de me rendre compte de l'absence de toute sensibilisation à l'IA dans le programme scolaire officiel car j'estime qu'il est capital de sensibiliser les élèves à ces problématiques, puisque l'usage de tels outils est de plus en plus répandu. Il est important que les jeunes générations soient conscientes des dangers que représente l'utilisation incontrôlée de tels outils, aussi bien d'un point de vue scolaire que d'un point de vue professionnel.
